% Forcing a registered self-appellation at the response opening — conference draft
% Compile: xelatex main.tex
\documentclass[10pt,twocolumn]{article}
\usepackage[a4paper,margin=2cm]{geometry}
\usepackage{booktabs}
\usepackage{graphicx}
\usepackage{amsmath}
\usepackage{url}
\usepackage{hyperref}

\title{Forcing a Registered Self-Appellation at the Response Opening\\Reduces LLM-Agent Sycophancy and Role Confusion:\\Mechanism Design and a Large-Scale Empirical Study Across Agent Platforms}
\author{Shiduki\\ (\href{https://github.com/teps3105}{teps3105})}
\date{November 2026 (draft)}

\begin{document}
\maketitle

\begin{abstract}
When large language model (LLM) agents converse with users, two recurring low-quality
response patterns appear: sycophancy (endorsing the user's stance without verification)
and role confusion (addressing tool output or non-human input as if it were a conversation
partner). We propose a zero-cost mechanism: force every agent response to begin with the
agent's registered self-appellation (in our deployment, a fictional agent name), enforced
by four layers of programmatic guards (pre-flight prompt injection, mid-turn interception,
end-of-turn blocking, and a first-character floor). The first-character time cost is
constant regardless of response length, yet it forces the model to establish its own
identity in the first generation step instead of continuing the user's utterance. We
measure five opening groups (self-appellation, addressing the user, first-person ``I'',
second-person ``you'', other) and six mechanically checkable quality indicators over
358{,}329 assistant responses across three agent platforms (Claude Code, Codex, ZCode).
Results: the self-appellation group (9{,}255 responses) has a zero sycophancy rate; the
association between second-person openings and sycophancy replicates across platforms
(40.9\% on Codex, 46.6\% on ZCode); role confusion is near 100\% in the addressee and
second-person groups. We also report post-deployment compliance, limitations, and threats
to validity.
\end{abstract}

\section{Introduction}
When LLMs act as interactive agents, two systematic response-quality patterns emerge.
The first is sycophancy: the model endorses the user's claim (``Right!'', ``You're
correct'') without performing any verification, which is known to stem from preference
learning in RLHF \cite{sharma2024towards}. The second is less studied: role confusion---
the agent treats tool execution results and system messages as a ``conversation partner'',
responding to machine output in the second person (``you''), indicating that the model
fails to locate its own identity relative to the listener during generation.

Existing mitigations (prompt guidance, RLHF fine-tuning) operate inside the model and are
hard to verify or maintain. We propose a mechanism that operates at the generation
interface: \textbf{force the first character of every response to be the agent's
registered self-appellation}. The intuition: opening with ``You're right'' versus opening
with ``\textit{〈agent name〉} has verified this'' forces different stance choices---the
former presupposes the user's stance, the latter forces the agent to start from its own
actions. The first-character time cost is constant regardless of response length.

Enforcement does not rely on model self-discipline but on four layers of programmatic
guards (\S\ref{sec:method}), which makes the effect measurable at scale: any violation
(user-addressing openings, English openings, sycophantic openings) is intercepted by
code. Our contributions:
\begin{enumerate}
  \item a deployable, verifiable, zero-generation-cost opening-appellation mechanism
        with a four-layer guard implementation;
  \item a grouped measurement framework with six mechanical indicators over 358{,}329
        real conversations across three platforms;
  \item empirical associations between opening groups and sycophancy/role confusion:
        zero sycophancy in the self-appellation group, and a cross-platform replication
        of the second-person/sycophancy association.
\end{enumerate}

\section{Related Work}
\textbf{LLM sycophancy.} Sharma et al.\ \cite{sharma2024towards} systematically study
sycophancy in RLHF models, showing that agreeing with the user is a rewarded behavior.
Mitigations include system prompts, data augmentation, and RLAIF. Our work differs: we
do not improve the model internally but force a different opening stance at the
generation interface, with programmatic enforcement. \textbf{Agent self-location.}
System prompts set personas (``you are assistant X''), but self-location during
generation is unconstrained. To our knowledge, forcing the first character of responses
to a registered self-appellation with programmatic verification has not been reported.
\textbf{Chinese response quality.} Simplified-character leakage and telegraphic phrasing
in agent settings have been reported without mechanical indicators; we integrate them
into the same measurement framework.

\section{Method}\label{sec:method}
\subsection{Mechanism}
The mechanism consists of four guard layers, all operating on the agent host's hook
interface (the extension point where the host runs user-defined programs around tool
calls and turn ends; supported by Claude Code, Codex, and ZCode):
\begin{enumerate}
  \item \textbf{Pre-flight injection} (UserPromptSubmit): inject a one-line rule hint
        when the user sends input;
  \item \textbf{Mid-turn interception} (PreToolUse): when narration addressed to the
        user is detected, deny the next tool call to force a rewrite;
  \item \textbf{End-of-turn blocking} (Stop): at turn end, inspect visible responses---
        six violation classes (user-addressing openings, English openings, sycophantic
        openings, jargon, simplified glyphs, typos) block the stop and require fixes;
  \item \textbf{First-character floor}: starting with the agent name forces the response
        to begin in Chinese---English cannot start (a side benefit).
\end{enumerate}
A verification gate exempts sycophancy blocks: any tool call in the same turn (fact
checking) allows agreement openings---the mechanism opposes unverified endorsement,
not agreement itself.

\subsection{Measurement Framework}
We scan the complete conversation records of three platforms: the SQLite conversation
databases of Claude Code and ZCode, and Codex's thread-history database---358{,}329
assistant responses in total. Each response's first non-empty text block is stripped of
Markdown prefixes and grouped into five groups: \textbf{self-appellation} (agent-name
opening), \textbf{addressee} (user-addressing), \textbf{first person} (``I''-opening),
\textbf{second person} (``you''-opening), \textbf{other}.

All six indicators are mechanically checkable (\S\ref{sec:metrics}). The classification
criteria and the online guards share the same code (the same module), avoiding drift
between measurement and enforcement criteria.

\section{Indicator Definitions}\label{sec:metrics}
\begin{table}[t]
\centering
\small
\begin{tabular}{p{3.2cm}p{4.6cm}p{1.6cm}}
\toprule
Indicator & Operational definition & Strength \\
\midrule
Sycophantic opening & First sentence hits an agreement-phrase list with zero tool calls
in the turn & Weak \\
Role confusion & Previous input is a tool result and the response opens addressing the
user & Strong \\
Mean length & Average response body characters & Descriptive \\
Simplified glyphs & Original glyphs that vanish under s2t conversion (curated list) & Strong \\
English opening & First non-empty line starts with an ASCII letter & Strong \\
Pronoun density & ``it/this/that'' occurrences per 1{,}000 characters & Weak \\
\bottomrule
\end{tabular}
\caption{Operational definitions and evidence strength. Weak indicators can be triggered
by harmless usage; strong indicators are almost always problems in agent contexts.}
\label{tab:metrics}
\end{table}

The core limitation of indicator strength: the sycophantic-opening criterion (opening
endorsement + zero verification) misses mid-sentence sycophancy and verified-but-
deferential tone---it measures the \textbf{surface layer} of sycophancy only. The role
confusion rate has virtually no false positives: addressing tool output as ``you'' has
no legitimate use in agent contexts.

\section{Results}
% Data: 2026-11-10, three-platform full corpus (Claude JSONL 37492, Codex thread_history
% 76170, ZCode sqlite 244667 responses)
\subsection{Opening-Group Distribution}
The three platforms total 358{,}329 assistant responses. The distribution is highly
uneven: self-appellation 9{,}255, addressee 9{,}390, ``I'' 14{,}445, ``you'' 1{,}342,
other 318{,}897 (mostly pre-deployment).

\begin{table}[t]
\centering
\small
\setlength{\tabcolsep}{3.5pt}
\begin{tabular}{llrrrrrr}
\toprule
Platform & Group & N & Chars & Syc.\% & Conf.\% & Pron.\% & Eng.\% \\
\midrule
Claude & Self-appel. & 5672 & 103 & 0.0 & 0.0 & 4.0 & 0.0 \\
Claude & Addressee & 3116 & 382 & 3.7 & 100 & 9.3 & 0.0 \\
Claude & I & 119 & 838 & 0.0 & 0.0 & 100 & 0.0 \\
Claude & You & 3 & 1309 & 0.0 & -- & 100 & 0.0 \\
Claude & Other & 28582 & 204 & 0.0 & 0.0 & 5.8 & 15.9 \\
\midrule
Codex & Self-appel. & 3569 & 226 & 0.0 & 0.0 & 0.3 & 0.0 \\
Codex & Addressee & 5647 & 180 & 1.1 & 100 & 0.8 & 0.0 \\
Codex & I & 9710 & 110 & 0.0 & 0.0 & 99.8 & 0.0 \\
Codex & You & 931 & 241 & 40.9 & 100 & 100 & 0.0 \\
Codex & Other & 56313 & 266 & 0.4 & 0.0 & 40.7 & 2.0 \\
\midrule
ZCode & Self-appel. & 14 & 924 & 0.0 & 0.0 & 21.4 & 0.0 \\
ZCode & Addressee & 6627 & 1699 & 0.6 & 100 & 73.3 & 0.4 \\
ZCode & I & 4616 & 1855 & 0.0 & 0.0 & 100 & 0.2 \\
ZCode & You & 408 & 588 & 46.6 & 100 & 100 & 0.0 \\
ZCode & Other & 233002 & 1584 & 0.0 & 0.0 & 25.2 & 37.3 \\
\bottomrule
\end{tabular}
\caption{Six indicators over five groups on three platforms. Conf.\ = addressing the
user right after a tool result (denominator: responses in the group whose previous input
was a tool result); the Claude ``you'' group has only 3 responses, insufficient for the
confusion column.}
\label{tab:main}
\end{table}

\subsection{Sycophancy and Opening Groups}
The second-person opening/sycophancy association replicates across platforms: the
``you'' group's sycophancy rate is 40.9\% on Codex (931 responses) and 46.6\% on ZCode
(408 responses). Addressee-group rates are markedly lower (Claude 3.7\%, Codex 1.1\%,
ZCode 0.6\%). The self-appellation group totals 9{,}255 responses across platforms with
a zero sycophancy rate. This matches the mechanism hypothesis: opening with the user's
appellation places the model in a ``responding to you'' position where agreement is the
natural continuation; opening with its own appellation forces the response to start
from the agent's own action.

\subsection{Role Confusion}
Role confusion is near 100\% in the addressee and ``you'' groups: Codex addressee
4{,}373/4{,}373, you-group 195/195; ZCode addressee 786/786, you-group 31/31. Both forms
of confusion (addressing the user and addressing tool output) share one opening habit---
when generation fails to locate the listener, the model says ``you'' to anyone. The
self-appellation and ``I'' groups show zero confusion.

\subsection{Response Length}
The length ratio between ``I''-openings and self-appellation openings is inconsistent
across platforms: 8.1$\times$ on Claude, 2.0$\times$ on ZCode, and an inverted
0.49$\times$ on Codex. Length is dominated by task type; we do not claim it as a
mechanism effect.

\section{Threats to Validity}
\textbf{Construct validity.} The sycophantic-opening rate is a surface-level proxy;
``no opening endorsement'' does not imply ``no content pandering''.
\textbf{Internal validity.} Groups are determined by the opening character, which the
mechanism itself enforces---post-deployment self-appellation responses are a direct
product of the mechanism, and their quality indicators are also shaped by other guards
(simplified glyphs, jargon, sentence length); attribution to the appellation alone is
not possible. \textbf{External validity.} The sample comes from a single user's (the
researcher's) agent workload; the conversation distribution differs from typical users.
The three platforms share the same global instructions and guards; between-platform
differences reflect model and interface differences, not prompt differences.
\textbf{Statistical power.} With large samples, proportional differences are almost
certainly significant; we therefore report effect sizes and raw distributions and do
not use p-values as primary evidence.

\section{Conclusion and Future Work}
Forcing a self-appellation opening is an interface-layer mechanism with zero cost and
programmatic verifiability. Current evidence shows systematic associations between
opening groups and sycophancy/role confusion. The mechanism's causal contribution
requires pre/post-deployment comparisons and ablations (injection only, blocking only,
all four layers), left as future work.

\begin{thebibliography}{1}
\bibitem{sharma2024towards}
M.~Sharma, M.~Tong, T.~Korbak, D.~Duvenaud, A.~Askell, S.~R.~Bowman, et al.
\newblock Towards understanding sycophancy in language models.
\newblock In \emph{ICLR}, 2024.
\end{thebibliography}

\end{document}
